% ================================================================== 13
\section{Referencias}
\label{sec:referencias}

\begingroup
\renewcommand{\section}[2]{}
\begin{thebibliography}{99}\small
\bibitem{clrs} T.~H. Cormen, C.~E. Leiserson, R.~L. Rivest y C.~Stein,
  \emph{Introduction to Algorithms}, 4.\textsuperscript{a} ed. MIT Press, 2022
  (cap.~31: algoritmos de teoría de números).
\bibitem{hac} A.~J. Menezes, P.~C. van Oorschot y S.~A. Vanstone,
  \emph{Handbook of Applied Cryptography}. CRC Press, 1996.
\bibitem{knuth} D.~E. Knuth, \emph{The Art of Computer Programming, Vol.~2:
  Seminumerical Algorithms}, 3.\textsuperscript{a} ed. Addison-Wesley, 1997.
\bibitem{rsa1978} R.~L. Rivest, A.~Shamir y L.~Adleman, «A method for obtaining
  digital signatures and public-key cryptosystems», \emph{Communications of the
  ACM}, vol.~21, n.º~2, pp.~120--126, 1978.
\bibitem{miller1976} G.~L. Miller, «Riemann's hypothesis and tests for
  primality», \emph{Journal of Computer and System Sciences}, vol.~13, n.º~3,
  pp.~300--317, 1976.
\bibitem{rabin1980} M.~O. Rabin, «Probabilistic algorithm for testing
  primality», \emph{Journal of Number Theory}, vol.~12, n.º~1, pp.~128--138, 1980.
\bibitem{pollard1975} J.~M. Pollard, «A Monte Carlo method for factorization»,
  \emph{BIT Numerical Mathematics}, vol.~15, n.º~3, pp.~331--334, 1975.
\bibitem{brent1980} R.~P. Brent, «An improved Monte Carlo factorization
  algorithm», \emph{BIT Numerical Mathematics}, vol.~20, n.º~2, pp.~176--184, 1980.
\bibitem{pomerance1980} C.~Pomerance, J.~L. Selfridge y S.~S. Wagstaff, «The
  pseudoprimes to $25\cdot10^9$», \emph{Mathematics of Computation}, vol.~35,
  n.º~151, pp.~1003--1026, 1980.
\bibitem{sorenson2017} J.~Sorenson y J.~Webster, «Strong pseudoprimes to twelve
  prime bases», \emph{Mathematics of Computation}, vol.~86, n.º~304,
  pp.~985--1003, 2017.
\bibitem{qc1982} J.-J. Quisquater y C.~Couvreur, «Fast decipherment algorithm
  for RSA public-key cryptosystem», \emph{Electronics Letters}, vol.~18, n.º~21,
  pp.~905--907, 1982.
\bibitem{rfc8017} K.~Moriarty, B.~Kaliski, J.~Jonsson y A.~Rusch,
  «PKCS \#1: RSA Cryptography Specifications Version 2.2», RFC~8017, IETF, 2016.
\bibitem{fips186} NIST, \emph{Digital Signature Standard (DSS)}, FIPS~186-5, 2023.
\bibitem{hypothesis} D.~R. MacIver, Z.~Hatfield-Dodds et~al., «Hypothesis: A new
  approach to property-based testing», \emph{Journal of Open Source Software},
  vol.~4, n.º~43, p.~1891, 2019.
\bibitem{pydocs} Python Software Foundation, documentación de los módulos
  \texttt{secrets}, \texttt{random}, \texttt{time} y \texttt{tracemalloc}.
  \url{https://docs.python.org/3/library/}
\bibitem{gmpy2} Documentación de gmpy2. \url{https://gmpy2.readthedocs.io}
\bibitem{pytest} Documentación de pytest. \url{https://docs.pytest.org}
\bibitem{nist80057} E.~Barker, \emph{Recommendation for Key Management: Part 1 -- General},
  NIST Special Publication 800-57 Part~1 Rev.~5, 2020.
\bibitem{aks2004} M.~Agrawal, N.~Kayal y N.~Saxena, «PRIMES is in P», \emph{Annals of Mathematics},
  vol.~160, n.º~2, pp.~781--793, 2004.
\bibitem{dlp1993} I.~Damgård, P.~Landrock y C.~Pomerance, «Average case error estimates for the
  strong probable prime test», \emph{Mathematics of Computation}, vol.~61, n.º~203, pp.~177--194, 1993.
\bibitem{bw1980} R.~Baillie y S.~S. Wagstaff, «Lucas pseudoprimes», \emph{Mathematics of
  Computation}, vol.~35, n.º~152, pp.~1391--1417, 1980.
\bibitem{lenstra1987} H.~W. Lenstra, «Factoring integers with elliptic curves», \emph{Annals of
  Mathematics}, vol.~126, n.º~3, pp.~649--673, 1987.
\bibitem{boudot2020} F.~Boudot, P.~Gaudry, A.~Guillevic, N.~Heninger, E.~Thomé y P.~Zimmermann,
  «Comparing the difficulty of factorization and discrete logarithm: a 240-digit experiment», en
  \emph{Advances in Cryptology -- CRYPTO 2020}, LNCS~12171, pp.~62--91, Springer, 2020.
\bibitem{progit} S.~Chacon y B.~Straub, \emph{Pro Git}, 2.\textsuperscript{a} ed. Apress, 2014.
  \url{https://git-scm.com/book}
\bibitem{pypa} Python Packaging Authority, \emph{Python Packaging User Guide}.
  \url{https://packaging.python.org}
\end{thebibliography}
\endgroup

\newpage
